% GENERATED FILE. Edit canonical lessons, then run npm run generate:books.
% canonical-source-hash: fnv1a64:b1de933cd39194a9
% canonical-lessons: IT-C435-riscaldamento, IT-C435-termosifone, IT-C435-condizionatore, IT-C435-rinfrescare, IT-C435-scaldabagno

\chapter{Caldo e freddo in camera}
\label{ch:it-caldo-e-freddo-in-camera}
\hlchaptermodality{435}

\begin{chapteropening}
\textbf{By the end of this chapter:} \emph{I can complain about the temperature: the heating, the radiator, the air conditioner, cooling down and the water heater.}

Complete the last lesson of chapter 435: I can complain about the temperature: the heating, the radiator, the air conditioner, cooling down and the water heater.
\end{chapteropening}

\section[il riscaldamento]{il riscaldamento — the heating}

\label{lesson:IT-C435-riscaldamento}



\begin{quote}
Before the new one: say the Italian for the drain, then the Italian for to unblock.
\end{quote}

\subsection*{You'll want to know: il riscaldamento}
\textbf{il riscaldamento} — \textquotedblleft{}the heating\textquotedblright{}.

\begin{quote}
Il riscaldamento è spento e fa freddo.
\end{quote}

\begin{quote}
\emph{The heating is off and it is cold.}
\end{quote}

\emph{Il riscaldamento globale} — global warming.

\begin{cousinweb}
\emph{riscaldamento} is built on \emph{scaldare} 'to warm', which is inherited from Late Latin \emph{excaldāre} (EtymDB), built on \emph{calidus} 'hot'.
\end{cousinweb}

\begin{grammarlens}[title={What you've built}]
One more word for telling, explaining and describing at B1.
\end{grammarlens}

\subsection*{Your turn}
\begin{itemize}
\raggedright
  \item \emph{Say it:} \emph{il riscaldamento}
  \item \emph{Say it:} \emph{il riscaldamento}, once more
  \item \emph{Say it:} a sentence of your own with \emph{il riscaldamento}
  \item \emph{Recall:} say the Italian for the drain, then the Italian for to unblock, then say \emph{il riscaldamento} again
\end{itemize}

\begin{tcolorbox}[breakable,colback=teal!4,colframe=teal!35!black,title={Before you move on}]
What does il riscaldamento mean? (\textquotedblleft{}The heating\textquotedblright{}.) Say it once more.
\end{tcolorbox}
\section[il termosifone]{il termosifone — the radiator}

\label{lesson:IT-C435-termosifone}



\begin{quote}
Before the new one: say the Italian for to unblock, then the Italian for the heating.
\end{quote}

\subsection*{You'll want to know: il termosifone}
\textbf{il termosifone} — \textquotedblleft{}the radiator\textquotedblright{}.

\begin{quote}
Metti i calzini sul termosifone.
\end{quote}

\begin{quote}
\emph{Put the socks on the radiator.}
\end{quote}

\emph{Accendere i termosifoni} — to switch the heating on.

\begin{cousinweb}
\emph{termosifone} is a learned formation from Greek \emph{thermós} 'hot' and \emph{sifone} 'siphon', from Greek \emph{síphōn} 'pipe, tube'.
\end{cousinweb}

\begin{grammarlens}[title={What you've built}]
One more word for telling, explaining and describing at B1.
\end{grammarlens}

\subsection*{Your turn}
\begin{itemize}
\raggedright
  \item \emph{Say it:} \emph{il termosifone}
  \item \emph{Say it:} \emph{il termosifone}, once more
  \item \emph{Say it:} a sentence of your own with \emph{il termosifone}
  \item \emph{Recall:} say the Italian for to unblock, then the Italian for the heating, then say \emph{il termosifone} again
\end{itemize}

\begin{tcolorbox}[breakable,colback=teal!4,colframe=teal!35!black,title={Before you move on}]
What does il termosifone mean? (\textquotedblleft{}The radiator\textquotedblright{}.) Say it once more.
\end{tcolorbox}
\section[il condizionatore]{il condizionatore — the air conditioner}

\label{lesson:IT-C435-condizionatore}



\begin{quote}
Before the new one: say the Italian for the heating, then the Italian for the radiator.
\end{quote}

\subsection*{You'll want to know: il condizionatore}
\textbf{il condizionatore} — \textquotedblleft{}the air conditioner\textquotedblright{}.

\begin{quote}
Il condizionatore fa molto rumore.
\end{quote}

\begin{quote}
\emph{The air conditioner is very noisy.}
\end{quote}

\emph{L'aria condizionata} is air conditioning.

\begin{cousinweb}
\emph{condizionatore} is built on \emph{condizionare} 'to condition', from \emph{condizione}, a learned borrowing from Latin \emph{condiciō} 'agreement, state' (Lewis \& Short).
\end{cousinweb}

\begin{grammarlens}[title={What you've built}]
One more word for telling, explaining and describing at B1.
\end{grammarlens}

\subsection*{Your turn}
\begin{itemize}
\raggedright
  \item \emph{Say it:} \emph{il condizionatore}
  \item \emph{Say it:} \emph{il condizionatore}, once more
  \item \emph{Say it:} a sentence of your own with \emph{il condizionatore}
  \item \emph{Recall:} say the Italian for the heating, then the Italian for the radiator, then say \emph{il condizionatore} again
\end{itemize}

\begin{tcolorbox}[breakable,colback=teal!4,colframe=teal!35!black,title={Before you move on}]
What does il condizionatore mean? (\textquotedblleft{}The air conditioner\textquotedblright{}.) Say it once more.
\end{tcolorbox}
\section[rinfrescare]{rinfrescare — to cool down, to freshen}

\label{lesson:IT-C435-rinfrescare}



\begin{quote}
Before the new one: say the Italian for the radiator, then the Italian for the air conditioner.
\end{quote}

\subsection*{You'll want to know: rinfrescare}
\textbf{rinfrescare} — \textquotedblleft{}to cool down, to freshen\textquotedblright{}.

\begin{quote}
La sera la stanza si rinfresca un po'.
\end{quote}

\begin{quote}
\emph{In the evening the room cools down a bit.}
\end{quote}

\emph{Rinfrescarsi} — to freshen up.

\begin{cousinweb}
\emph{rinfrescare} is built on \emph{fresco} 'cool, fresh', which was borrowed through Medieval Latin \emph{friscus} from a Germanic word (EtymDB).
\end{cousinweb}

\begin{grammarlens}[title={What you've built}]
One more word for telling, explaining and describing at B1.
\end{grammarlens}

\subsection*{Your turn}
\begin{itemize}
\raggedright
  \item \emph{Say it:} \emph{rinfrescare}
  \item \emph{Say it:} \emph{rinfrescare}, once more
  \item \emph{Say it:} a sentence of your own with \emph{rinfrescare}
  \item \emph{Recall:} say the Italian for the radiator, then the Italian for the air conditioner, then say \emph{rinfrescare} again
\end{itemize}

\begin{tcolorbox}[breakable,colback=teal!4,colframe=teal!35!black,title={Before you move on}]
What does rinfrescare mean? (\textquotedblleft{}To cool down, to freshen\textquotedblright{}.) Say it once more.
\end{tcolorbox}
\section[lo scaldabagno]{lo scaldabagno — the water heater}

\label{lesson:IT-C435-scaldabagno}



\begin{quote}
Before the new one: say the Italian for the air conditioner, then the Italian for to cool down.
\end{quote}

\subsection*{You'll want to know: lo scaldabagno}
\textbf{lo scaldabagno} — \textquotedblleft{}the water heater\textquotedblright{}.

\begin{quote}
Lo scaldabagno è rotto: niente acqua calda.
\end{quote}

\begin{quote}
\emph{The water heater is broken: no hot water.}
\end{quote}

In speech also \emph{il boiler}. Plural unchanged.

\begin{cousinweb}
\emph{scaldabagno} is built from \emph{scaldare} 'to warm', inherited from Late Latin \emph{excaldāre}, and \emph{bagno} 'bath', inherited from Latin \emph{balneum} (EtymDB).
\end{cousinweb}

\begin{grammarlens}[title={What you've built}]
One more word for telling, explaining and describing at B1.
\end{grammarlens}

\subsection*{Your turn}
\begin{itemize}
\raggedright
  \item \emph{Say it:} \emph{lo scaldabagno}
  \item \emph{Say it:} \emph{lo scaldabagno}, once more
  \item \emph{Say it:} a sentence of your own with \emph{lo scaldabagno}
  \item \emph{Recall:} say the Italian for the air conditioner, then the Italian for to cool down, then say \emph{lo scaldabagno} again
\end{itemize}

\begin{tcolorbox}[breakable,colback=teal!4,colframe=teal!35!black,title={Before you move on}]
What does lo scaldabagno mean? (\textquotedblleft{}The water heater\textquotedblright{}.) Say it once more.
\end{tcolorbox}
